\documentclass[aspectratio=169,11pt]{beamer}

% Самостоятельный исходник: pdfLaTeX, без внешних рисунков и файлов.
\usepackage{iftex}
\ifPDFTeX
  \usepackage[T2A]{fontenc}
  \usepackage[utf8]{inputenc}
\else
  \usepackage{fontspec}
  \setmainfont{cmunrm.otf}[BoldFont=cmunbx.otf,ItalicFont=cmunti.otf,BoldItalicFont=cmunbi.otf]
  \setsansfont{cmunss.otf}[BoldFont=cmunsx.otf,ItalicFont=cmunsi.otf]
  \setmonofont{cmuntt.otf}
\fi
\usepackage[russian]{babel}
\usepackage{amsmath,amssymb,amsfonts,array}
\usepackage{tikz}
\usetikzlibrary{arrows.meta}

\usetheme{Madrid}
\usefonttheme{serif}
\setbeamercolor{background canvas}{bg=white}
\setbeamercolor{structure}{fg=blue!70!black}
\setbeamertemplate{navigation symbols}{}
\setbeamertemplate{footline}{%
  \leavevmode\hbox{%
  \begin{beamercolorbox}[wd=.28\paperwidth,ht=2.6ex,dp=1.2ex,center]{author in head/foot}%
  \insertshortauthor{} (\insertshortinstitute)\end{beamercolorbox}%
  \begin{beamercolorbox}[wd=.50\paperwidth,ht=2.6ex,dp=1.2ex,center]{title in head/foot}%
  Обыкновенные дифференциальные уравнения\end{beamercolorbox}%
  \begin{beamercolorbox}[wd=.22\paperwidth,ht=2.6ex,dp=1.2ex,center]{date in head/foot}%
  Лекция 1 \quad \insertframenumber/\inserttotalframenumber\end{beamercolorbox}}}
\newcommand{\R}{\mathbb R}
\newcommand{\dd}{\,d}
\title[Понятие ОДУ]{Обыкновенные дифференциальные уравнения}
\subtitle{Лекция 1. Понятие дифференциального уравнения}
% Лектор и место проведения курса.
\author[Ховрычев А. А.]{Ховрычев Андрей Александрович}
\institute[УрФУ]{ФИИТ, УрФУ}
\date{}

\begin{document}

% 1
\begin{frame}
  \titlepage
  \centering От закона изменения к движению
\end{frame}

\section{От модели к уравнению}
% 2
\begin{frame}{Что мы хотим узнать о движении?}
  Пусть $x(t)$ --- температура, численность популяции или координата.
  \begin{block}{От локального закона к прогнозу}
    Зная, \emph{как быстро меняется} величина сейчас, можно ли определить,
    \emph{какой она станет} через некоторое время?
  \end{block}
  К концу лекции мы научимся:
  \begin{itemize}
    \item записывать простую модель через производную;
    \item проверять решение и учитывать начальные данные;
    \item читать поле направлений;
    \item различать существование, единственность и время жизни решения.
  \end{itemize}
\end{frame}

% 3
\begin{frame}{Модель: как остывает чашка кофе?}
  Обозначим через $T(t)$ температуру кофе, через $T_{\text{ср}}$ ---
  постоянную температуру окружающей среды.
  \begin{block}{Модельное предположение}
    Скорость изменения температуры пропорциональна разности температур:
    \[
      \boxed{T'(t)=-k\bigl(T(t)-T_{\text{ср}}\bigr)},\qquad k>0.
    \]
  \end{block}
  \begin{itemize}
    \item Если $T>T_{\text{ср}}$, то $T'<0$: кофе остывает.
    \item Если $T<T_{\text{ср}}$, то $T'>0$: содержимое нагревается.
    \item Размерность $k$ --- обратное время, например $\text{мин}^{-1}$.
  \end{itemize}
  Это приближённая модель: среда и коэффициент теплообмена считаются постоянными.
\end{frame}

% 4
\begin{frame}{Одного закона недостаточно}
  Введём отклонение от температуры среды:
  \[
    x(t)=T(t)-T_{\text{ср}}.\qquad \text{Тогда } x'=-kx.
  \]
  При любом $C\in\R$ функция $x(t)=Ce^{-kt}$ удовлетворяет уравнению:
  \[
    x'(t)=-kCe^{-kt}=-kx(t).
  \]
  \begin{alertblock}{Чего не хватает для прогноза?}
    Закон охлаждения одинаков для чашек с разной начальной температурой.
    Чтобы выбрать движение, нужно задать, например, $x(0)=x_0$.
  \end{alertblock}
  Тогда $C=x_0$ и $T(t)=T_{\text{ср}}+(T(0)-T_{\text{ср}})e^{-kt}$.
\end{frame}

% 5
\begin{frame}{Разные модели --- общий язык}
  \renewcommand{\arraystretch}{1.6}
  \begin{tabular}{>{\raggedright\arraybackslash}p{.29\textwidth}>{\raggedright\arraybackslash}p{.32\textwidth}>{\raggedright\arraybackslash}p{.29\textwidth}}
    \textbf{Процесс} & \textbf{Уравнение} & \textbf{Смысл закона} \\
    Рост популяции & $N'=rN$, $r>0$ & Скорость пропорциональна численности \\
    Свободное падение & $h''=-g$, $g>0$ & Ускорение постоянно; ось направлена вверх \\
    Груз на пружине & $mq''=-\kappa q$ & Возвращающая сила пропорциональна смещению
  \end{tabular}
  \vspace{.5em}
  \begin{block}{Общий принцип}
    Сначала выбираем величины и допущения, затем выражаем закон через производные.
  \end{block}
\end{frame}

\section{Уравнение и решение}
% 6
\begin{frame}{Что такое обыкновенное дифференциальное уравнение?}
  \begin{block}{Определение}
    \alert{Обыкновенное дифференциальное уравнение (ОДУ)} связывает
    независимую переменную, неизвестную функцию этой переменной
    и её производные:
    \[
      F\bigl(t,x,x',\ldots,x^{(n)}\bigr)=0.
    \]
  \end{block}
  \alert{Порядок} --- наивысший порядок производной, действительно входящей в уравнение.
  \[
    x'=t-x\quad\text{--- первый порядок},\qquad
    x''+x=0\quad\text{--- второй порядок}.
  \]
  Неизвестная может быть вектором: тогда получаем систему ОДУ.
\end{frame}

% 7
\begin{frame}{Почему «обыкновенное»?}
  В ОДУ неизвестная функция зависит от \alert{одной} независимой переменной:
  \[
    x=x(t),\qquad x' = \frac{dx}{dt}=\dot x.
  \]
  Буква $t$ удобна для времени, но переменная может означать и координату.
  \begin{block}{Другой класс задач}
    Если $u=u(t,s)$ зависит от нескольких переменных, возникают
    уравнения с частными производными, например
    \[
      \frac{\partial u}{\partial t}
      =a\frac{\partial^2u}{\partial s^2},\qquad a>0.
    \]
  \end{block}
  Система из нескольких функций одного $t$ всё ещё является системой ОДУ.
\end{frame}

% 8
\begin{frame}{Нормальная форма и область определения}
  Основной объект в начале курса:
  \[
    \boxed{x'=f(t,x)},\qquad (t,x)\in D\subset\R^2,
  \]
  где $D$ --- открытая область, на которой задана правая часть.
  \begin{exampleblock}{Область --- часть постановки}
    Для $x'=1/x$ естественно взять $D=\{(t,x):x\ne0\}$.
    Начальная точка с $x_0=0$ в эту область не входит.
  \end{exampleblock}
  \begin{alertblock}{Не всякое ОДУ имеет одну глобальную нормальную форму}
    Уравнение $(x')^2=x$ при $x>0$ допускает две ветви:
    $x'=\sqrt{x}$ и $x'=-\sqrt{x}$.
  \end{alertblock}
\end{frame}

% 9
\begin{frame}{Что называется решением?}
  \begin{block}{Определение для $x'=f(t,x)$}
    Функция $\varphi:I\to\R$, где $I$ --- открытый интервал,
    является \alert{решением}, если:
    \begin{enumerate}
      \item $\varphi$ непрерывно дифференцируема на $I$;
      \item $(t,\varphi(t))\in D$ при каждом $t\in I$;
      \item $\varphi'(t)=f(t,\varphi(t))$ при каждом $t\in I$.
    \end{enumerate}
  \end{block}
  Решение задаётся \alert{функцией вместе с интервалом}.
  Для уравнения порядка $n$ аналогично требуется $\varphi\in C^n(I)$.
  \medskip

  На замкнутых временных отрезках будем рассматривать ограничения таких решений.
\end{frame}

% 10
\begin{frame}{Проверить решение проще, чем найти его}
  Рассмотрим $x'=x^2$ и $\varphi(t)=1/(1-t)$ на $I=(-\infty,1)$.
  \[
    \varphi'(t)=\frac{1}{(1-t)^2}=\varphi(t)^2.
  \]
  Значит, это решение. Кроме того, $\varphi(0)=1$.
  \begin{exampleblock}{Ещё одно решение}
    $\psi(t)\equiv0$ на $\R$: действительно, $\psi'=0=\psi^2$.
  \end{exampleblock}
  \begin{alertblock}{Интервал нельзя опускать}
    Та же формула $1/(1-t)$ задаёт решение на $(1,+\infty)$,
    но это не решение задачи с начальным моментом $t=0$.
  \end{alertblock}
\end{frame}

% 11
\begin{frame}{Все решения уравнения $x'=kx$}
  Пусть $k\in\R$ постоянно. Подстановка показывает, что $Ce^{kt}$ --- решение.
  Докажем, что других решений на интервале нет.
  \[
    \frac{d}{dt}\bigl(e^{-kt}x(t)\bigr)
    =e^{-kt}\bigl(x'(t)-kx(t)\bigr)=0.
  \]
  Функция с нулевой производной постоянна на интервале, поэтому
  \[
    e^{-kt}x(t)=C,\qquad \boxed{x(t)=Ce^{kt}}.
  \]
  \begin{block}{Два разных шага}
    Проверить найденные функции --- ещё не значит доказать,
    что найдены \emph{все} решения.
  \end{block}
\end{frame}

% 12
\begin{frame}{Задача Коши}
  \begin{block}{Определение}
    Найти решение уравнения, принимающее заданное значение
    в заданный начальный момент:
    \[
      \boxed{x'=f(t,x),\qquad x(t_0)=x_0},\qquad (t_0,x_0)\in D.
    \]
    Интервал решения должен содержать $t_0$.
  \end{block}
  Для $x'=kx$ получаем
  \[
    x_0=Ce^{kt_0},\qquad
    \boxed{x(t)=x_0e^{k(t-t_0)}}.
  \]
  В этом примере решение существует на всей $\R$ и единственно.
\end{frame}

% 13
\begin{frame}{Начальные данные выбирают кривую}
  \begin{columns}[c]
    \begin{column}{.55\textwidth}
      \centering
      \begin{tikzpicture}[x=1.35cm,y=.75cm,>=Stealth]
        \draw[->] (-.2,0)--(3.9,0) node[right] {$t$};
        \draw[->] (0,-2.3)--(0,3.5) node[above] {$x$};
        \foreach \c in {-2,-1,1,3}{
          \draw[blue!35,thick,domain=0:3.5,samples=50] plot (\x,{\c*exp(-\x)});
        }
        \draw[blue!80!black,very thick,domain=0:3.5,samples=50] plot (\x,{2*exp(-\x)});
        \fill[red!75!black] (0,2) circle (2pt);
        \node[right] at (.1,2.35) {$x(0)=2$};
        \node[below left] at (0,0) {$0$};
        \node at (2,3) {$x'=-x$};
      \end{tikzpicture}
    \end{column}
    \begin{column}{.43\textwidth}
      Семейство решений:
      \[x(t)=Ce^{-t}.
      \]
      Начальное условие $x(0)=2$ выбирает
      \[x(t)=2e^{-t}.
      \]
      Постоянное решение $x=0$ также входит в семейство.
    \end{column}
  \end{columns}
\end{frame}

% 14
\begin{frame}{Сколько начальных данных нужно?}
  Для нормального уравнения второго порядка
  \[
    x''=F(t,x,x')
  \]
  в задаче Коши задают $x(t_0)=x_0$ и $x'(t_0)=v_0$.
  Например, для свободного падения:
  \[
    h''=-g,\qquad h(0)=h_0,\qquad h'(0)=v_0,
  \]
  \[
    h(t)=h_0+v_0t-\frac{gt^2}{2}.
  \]
  \begin{block}{Порядок и данные}
    Для нормального уравнения порядка $n$ задают значения
    $x,x',\ldots,x^{(n-1)}$ в один момент.
    Существование и единственность требуют условий на правую часть.
  \end{block}
\end{frame}

\section{Геометрия решений}
% 15
\begin{frame}{Производная задаёт направление}
  В точке $(t,x)$ уравнение $x'=f(t,x)$ задаёт наклон касательной
  к графику решения.
  \[
    \text{Направляющий вектор касательной:}\qquad (1,f(t,x)).
  \]
  Совокупность таких направлений называется \alert{полем направлений}.
  \begin{block}{Интегральная кривая}
    График $x=\varphi(t)$ решения всюду касается заданного поля.
  \end{block}
  Для $x'=t-x$:
  \begin{itemize}
    \item ниже прямой $x=t$ наклоны положительны;
    \item выше неё --- отрицательны;
    \item на прямой $x=t$ касательные горизонтальны.
  \end{itemize}
\end{frame}

% 16
\begin{frame}{Поле направлений: $x'=t-x$}
  \begin{columns}[c]
    \begin{column}{.60\textwidth}
      \centering
      \begin{tikzpicture}[x=1.35cm,y=.85cm,>=Stealth]
        \draw[->] (-.2,0)--(3.65,0) node[right] {$t$};
        \draw[->] (0,-1.35)--(0,3.5) node[above] {$x$};
        \foreach \t in {0,.4,...,3.2}{
          \foreach \y in {-1,-.5,...,3}{
            \pgfmathsetmacro{\s}{\t-\y}
            \pgfmathsetmacro{\a}{.12/sqrt(1+\s*\s)}
            \draw[gray!70] ({\t-\a},{\y-\a*\s})--({\t+\a},{\y+\a*\s});
          }
        }
        \draw[dashed,orange!85!black] (0,0)--(3.2,3.2);
        \draw[blue!75!black,very thick,domain=0:3.2,samples=70]
          plot (\x,{\x-1+exp(-\x)});
        \draw[red!70!black,thick,domain=0:3.2,samples=70]
          plot (\x,{\x-1+3*exp(-\x)});
        \fill[blue!75!black] (0,0) circle (2pt);
        \fill[red!70!black] (0,2) circle (2pt);
        \node[orange!85!black,right] at (2.6,3.2) {$x=t$};
      \end{tikzpicture}
    \end{column}
    \begin{column}{.38\textwidth}
      Синяя кривая: $x(0)=0$.
      \medskip

      Красная кривая: $x(0)=2$.
      \medskip

      На пунктирной прямой наклон поля равен нулю.
      \begin{alertblock}{Вопрос}
        Является ли сама прямая $x=t$ решением?
      \end{alertblock}
    \end{column}
  \end{columns}
\end{frame}

% 17
\begin{frame}{Нулевой наклон не означает решение}
  Для предполагаемого решения $x(t)=t$ имеем
  \[
    x'(t)=1,\qquad t-x(t)=0.
  \]
  Левая и правая части не равны. Значит, $x=t$ \alert{не является решением}.
  \begin{block}{Что показывает линия $f(t,x)=0$?}
    На ней у проходящих через неё решений горизонтальные касательные.
    Самой этой линии не обязательно быть интегральной кривой.
  \end{block}
  Для сравнения, постоянная функция $x(t)=a$ является решением
  тогда и только тогда, когда $f(t,a)=0$ на всём рассматриваемом интервале.
\end{frame}

% 18
\begin{frame}{Автономные уравнения и равновесия}
  \begin{block}{Определение}
    Уравнение $x'=h(x)$ называется \alert{автономным}:
    правая часть не зависит явно от времени.
    Точка $a$ с $h(a)=0$ задаёт \alert{равновесие} $x(t)\equiv a$.
  \end{block}
  Для $x'=-x$ равновесие $x=0$; при $x>0$ величина убывает,
  при $x<0$ --- возрастает.
  \begin{center}
    \begin{tikzpicture}[>=Stealth,x=1.25cm]
      \draw[->] (-3,0)--(3,0) node[right] {$x$};
      \fill (0,0) circle (2pt) node[below] {$0$};
      \draw[very thick,blue!75!black,->] (-2.5,0)--(-.7,0);
      \draw[very thick,blue!75!black,->] (2.5,0)--(.7,0);
    \end{tikzpicture}
  \end{center}
  Это \alert{фазовая прямая}: она показывает направление изменения состояния.
  Из формулы $x(t)=x_0e^{-t}$ видно, что $x(t)\to0$ при $t\to+\infty$.
\end{frame}

\section{Три вопроса о задаче Коши}
% 19
\begin{frame}{Что нужно обосновать, кроме формулы?}
  Для задачи $x'=f(t,x)$, $x(t_0)=x_0$ зададим три вопроса.
  \begin{enumerate}
    \item \alert{Существование.} Есть ли решение хотя бы вблизи $t_0$?
    \item \alert{Единственность.} Может ли через начальную точку
      проходить несколько решений?
    \item \alert{Продолжение.} На каком интервале можно следить за решением?
  \end{enumerate}
  \begin{block}{Позже появится четвёртый вопрос}
    Насколько изменится решение при небольшой ошибке в начальных данных?
  \end{block}
  Ответы определяются свойствами уравнения, а не только числом заданных условий.
\end{frame}

% 20
\begin{frame}{Гладкое уравнение, но конечное время жизни}
  Рассмотрим
  \[
    x'=x^2,\qquad x(0)=1.
  \]
  Мы уже проверили решение
  \[
    x(t)=\frac{1}{1-t},\qquad -\infty<t<1.
  \]
  При $t\to1-0$ значение $x(t)\to+\infty$.
  \begin{alertblock}{Продолжить через $t=1$ нельзя}
    Решение с конечным значением в $t=1$ должно быть непрерывным там.
    У данной функции конечного предела нет.
  \end{alertblock}
  Правая часть $f(t,x)=x^2$ при этом гладкая на всей плоскости.
\end{frame}

% 21
\begin{frame}{Максимальный интервал решения}
  \begin{block}{Определение}
    Решение максимально, если его нельзя продолжить как решение
    того же уравнения на больший интервал.
  \end{block}
  Для $x(t)=1/(1-t)$ с $x(0)=1$:
  \begin{itemize}
    \item $(-1/2,1/2)$ --- допустимый, но не максимальный интервал;
    \item $(-\infty,1)$ --- максимальный интервал;
    \item $\R\setminus\{1\}$ --- не интервал и не область одного решения
      в принятом определении.
  \end{itemize}
  \begin{block}{Математика и модель}
    Математическое решение падения продолжается после удара о землю,
    но исходная физическая модель с этого момента требует изменения.
  \end{block}
\end{frame}

% 22
\begin{frame}{Начальные данные не всегда выбирают одно решение}
  Рассмотрим непрерывную правую часть и начальное условие
  \[
    x'=2\sqrt{|x|},\qquad x(0)=0.
  \]
  Для любого $a\ge0$ определим на всей $\R$
  \[
    x_a(t)=
    \begin{cases}
      0,&t\le a,\\
      (t-a)^2,&t>a.
    \end{cases}
  \]
  При $t<a$ обе части уравнения равны нулю; при $t>a$ --- $2(t-a)$.
  В точке $a$ обе односторонние производные равны нулю, и $x_a\in C^1(\R)$.
  \begin{alertblock}{Неединственность}
    Все $x_a$ решают одну задачу Коши; решение $x\equiv0$ также допустимо.
  \end{alertblock}
\end{frame}

% 23
\begin{frame}{Можно оставаться на месте, а затем начать движение}
  \begin{columns}[c]
    \begin{column}{.56\textwidth}
      \centering
      \begin{tikzpicture}[x=1.15cm,y=.9cm,>=Stealth]
        \draw[->] (-.65,0)--(4.1,0) node[right] {$t$};
        \draw[->] (0,-.25)--(0,3.8) node[above] {$x$};
        \foreach \a/\col in {0/blue,1/red,2/teal}{
          \draw[\col!75!black,very thick,domain=\a:{\a+1.8},samples=50]
            plot (\x,{(\x-\a)^2});
          \fill[\col!75!black] (\a,0) circle (2pt);
          \node[below] at (\a,0) {$\a$};
        }
        \draw[very thick,gray] (-.5,0)--(3.8,0);
        \node[blue!75!black] at (.6,2.4) {$a=0$};
        \node[red!75!black] at (1.8,2.5) {$a=1$};
        \node[teal!75!black] at (3,2.7) {$a=2$};
      \end{tikzpicture}
    \end{column}
    \begin{column}{.42\textwidth}
      До момента $a$ решение равно нулю, после --- растёт как $(t-a)^2$.
      \medskip

      Поле направлений само по себе не гарантирует единственность.
      \begin{block}{Что обсудим в лекции 3?}
        Какие дополнительные условия на $f$ исключают такое ветвление?
      \end{block}
    \end{column}
  \end{columns}
\end{frame}

% 24
\begin{frame}{Пересекаются ли интегральные кривые?}
  Если для каждой точки области выполнены условия \alert{локальной единственности},
  два решения, встретившись в точке $(t_*,x_*)$, совпадают
  на общем интервале определения, содержащем $t_*$.
  \begin{block}{Почему?}
    В точке встречи они решают одну задачу Коши. Локальная единственность
    даёт совпадение поблизости, а затем --- по всему общему интервалу.
  \end{block}
  Уравнение $x'=2\sqrt{|x|}$ показывает, что без условий единственности
  решения могут совпадать некоторое время, а затем расходиться.
  \medskip

  Сейчас это условный вывод. Достаточные условия единственности
  сформулируем и докажем в лекции 3.
\end{frame}

\section{Как мы будем изучать ОДУ}
% 25
\begin{frame}{Три подхода к одному уравнению}
  \begin{enumerate}
    \item \alert{Аналитический:} найти формулу и проверить её.
      \[x'=-kx\quad\Longrightarrow\quad x(t)=x_0e^{-kt}.\]
    \item \alert{Геометрический:} выяснить, растёт ли решение,
      куда стремится, где находятся равновесия.
    \item \alert{Численный:} вычислить приближённые значения
      в последовательные моменты времени.
  \end{enumerate}
  \begin{block}{Вопрос определяет инструмент}
    Для прогноза и анализа не всегда нужна явная формула.
    Но приближённый расчёт требует контроля ошибки.
  \end{block}
\end{frame}

% 26
\begin{frame}{Первый численный шаг: движение по касательной}
  Из определения производной для малого $h$:
  \[
    x(t_0+h)=x_0+h f(t_0,x_0)+o(h).
  \]
  Отбрасывая остаток, получаем приближение
  \[
    \boxed{x_1=x_0+h f(t_0,x_0)}.
  \]
  Например, при $x'=-x$, $x(0)=1$, $h=0.1$:
  \[
    x_1=0.9,\qquad x(0.1)=e^{-0.1}\approx0.90484.
  \]
  \begin{alertblock}{Приближение не является точным решением}
    При повторении шагов ошибки накапливаются. Метод Эйлера,
    выбор шага и его ограничения --- тема следующей лекции.
  \end{alertblock}
\end{frame}

% 27
\begin{frame}{Проверим понимание}
  \begin{enumerate}
    \item Уравнение $x'=x^2$ имеет решение $x=0$?
    \item Можно ли назвать $1/(1-t)$ решением на всей $\R$?
    \item Для $x''+x=0$ достаточно ли условия $x(0)=0$,
      чтобы выбрать одно решение?
    \item Если $f(t,x)$ непрерывна, обязана ли задача Коши
      иметь единственное решение?
    \item Является ли линия нулевого наклона $x=t$ решением $x'=t-x$?
  \end{enumerate}
  \medskip
  Для каждого ответа приведите проверку или контрпример.
\end{frame}

% 28
\begin{frame}{Что нужно унести с первой лекции}
  \begin{itemize}
    \item ОДУ описывает связь состояния и скорости его изменения.
    \item Решение --- функция на интервале, удовлетворяющая уравнению
      в каждой точке этого интервала.
    \item Задача Коши добавляет начальные данные к закону изменения.
    \item Поле направлений связывает уравнение с геометрией графиков.
    \item Существование, единственность и продолжимость --- разные вопросы.
  \end{itemize}
  \begin{block}{Следующий шаг}
    Научимся решать простейшие уравнения, исследовать автономную динамику
    на прямой и строить приближения методом Эйлера.
  \end{block}
\end{frame}

% 29
\begin{frame}{Упражнения после лекции}
  \small
  \begin{enumerate}
    \item Проверьте, что $x(t)=t^2+C$ решает $x'=2t$.
      Найдите решение с $x(1)=0$ и докажите, что других нет.
    \item Для $x'=x^2$, $x(0)=-1$ проверьте формулу
      $x(t)=-1/(t+1)$ и найдите её максимальный интервал, содержащий $0$.
    \item При $T_{\text{ср}}=20$, $T(0)=80$, $T(10)=50$
      найдите $k$ в модели охлаждения и предскажите $T(20)$.
      Время измеряется в минутах, температура --- в градусах Цельсия.
    \item Найдите $A,B$ так, чтобы $x(t)=A\cos t+B\sin t$
      решала $x''+x=0$, $x(0)=2$, $x'(0)=-1$.
    \item Нарисуйте несколько направлений для $x'=t-x$.
      Покажите, что $x=t-1+Ce^{-t}$ --- семейство решений.
    \item Объясните, почему $x_a$ со слайда 22 дифференцируема в $t=a$.
      Как этот пример опровергает утверждение «одно условие --- одно решение»?
  \end{enumerate}
\end{frame}

% 30
\begin{frame}{Материалы для самостоятельной работы}
  \begin{itemize}
    \item \href{https://ocw.mit.edu/courses/18-03sc-differential-equations-fall-2011/}{MIT OpenCourseWare: 18.03SC Differential Equations.}
      Раздел Unit I: базовые уравнения, геометрические и численные методы.
    \item \href{https://ocw.mit.edu/courses/18-03-differential-equations-spring-2010/pages/lecture-notes/}{MIT 18.03: Lecture Notes.}
      Материалы по уравнениям первого порядка и моделированию.
    \item План нашего курса: \texttt{plan.md}.
    \item Вопросы к аудитории и разбор упражнений:
      \texttt{lecture1-script.md}.
  \end{itemize}
  \begin{block}{Как работать с материалом}
    Перед поиском формулы назовите неизвестную функцию, область,
    начальные данные и вопрос, на который должна ответить модель.
  \end{block}
\end{frame}

\end{document}
